\documentclass[../../main.tex]{subfiles}% 注意这里的文件路径不能用 ./main.tex ,否则用latexmk编译子文件会报错
\graphicspath{{\subfix{./image/}}{\subfix{../../image/}}} % 指定图片引用路径,后续可以直接使用图片文件名
% 如果图片存放在上级目录,则使用 ../../image/ 作为备用路径

% 例如:
% \begin{figure}[H]
% \centering
% \includegraphics[width=0.3\textwidth]{图.png}
% \caption{}
% \label{figure:图}
% \end{figure}
% 注意:上述\label{}一定要放在\caption{}之后,否则引用图片序号会只会显示??.

\begin{document}

\section{极值原理和最大模估计}

\subsection{极值原理}

\begin{theorem}\label{theorem:PDE-极值原理1}
设$\Omega$ 是 $\mathbb{R}^n$ 上的有界开集,$c(x)\geqslant 0$，$f(x)<0$。如果 $u\in C^2(\Omega)\cap C(\overline{\Omega})$ 满足方程 \eqref{eq:general-elliptic}
\begin{align}
\mathcal{L}u \triangleq -\Delta u + c(x)u = f(x), \quad x\in\Omega,\label{eq:general-elliptic}
\end{align}
则 $u(x)$ 不能在 $\Omega$ 上达到它在 $\overline{\Omega}$ 上的非负最大值，即 $u(x)$ 只能在 $\partial\Omega$ 上达到它的非负最大值。
\end{theorem}
\begin{proof}
用反证法。如果 $u(x)$ 在点 $x_0\in\Omega$ 达到非负最大值，即
\begin{align*}
u(x_0) = \max_{x\in\Omega} u(x) \geqslant 0,
\end{align*}
则由\hyperref[Basis of Analytics-theorem:多元函数极值原理]{多元函数极值原理}知，$u(x)$ 在 $x_0$ 的梯度向量 $Du(x_0)=0$ 和 Hessian 矩阵 $D^2u(x_0)$ 是半负定的。对 Hessian 矩阵 $D^2u(x_0)$ 求迹得到
\begin{align*}
\Delta u(x_0) = \operatorname{tr}(D^2u(x_0)) \leqslant 0,
\end{align*}
因而
\begin{align*}
\mathcal{L}u(x_0) = -\Delta u(x_0) + c(x_0)u(x_0) = f(x_0) \geqslant 0.
\end{align*}
这与定理的假设 $f(x_0)<0$ 矛盾，因此 $u(x)$ 不能在 $\Omega$ 上达到它的非负最大值.
\end{proof}

\begin{theorem}\label{theorem:Lu方程极值原理}
设$\Omega$ 是 $\mathbb{R}^n$ 上的有界开集,$c(x)\geqslant 0$，$f(x)\leqslant 0$。如果 $u\in C^2(\Omega)\cap C(\overline{\Omega})$ 满足方程
\begin{align*}
\mathcal{L}u = -\Delta u + c(x)u = f(x), \quad x\in\Omega,
\end{align*}
则 $u(x)$ 必在 $\partial\Omega$ 上达到它在 $\overline{\Omega}$ 上的非负最大值，即
\begin{align*}
\max_{x\in\overline{\Omega}} u(x) \leqslant \max_{x\in\partial\Omega} u^+(x),
\end{align*}
其中 $u^+(x)=\max\{u,0\}$。
\end{theorem}
\begin{remark}
如果 $u(x)$ 在 $\overline{\Omega}$ 上的最大值是负数，那么上述\refthe{theorem:Lu方程极值原理}并没有告诉我们任何结论。
\end{remark}
\begin{proof}
不妨设原点 $0\in\Omega$。记 $d$ 为 $\Omega$ 的直径。对于任意 $\varepsilon>0$，我们构造辅助函数
\begin{align*}
w(x) = u(x) - \varepsilon(d^2 - |x|^2).
\end{align*}
显然
\begin{align*}
w(x) \leqslant u(x) \leqslant w(x) + \varepsilon d^2.
\end{align*}
计算得到
\begin{align*}
\mathcal{L}w &= \mathcal{L}u - 2n\varepsilon - c(x)\varepsilon(d^2 - |x|^2) \leqslant f - 2n\varepsilon < 0.
\end{align*}
由\refthe{theorem:PDE-极值原理1}知，$w$ 的非负最大值只能在 $\partial\Omega$ 上达到，因此
\begin{align*}
\max_{x\in\overline{\Omega}} w(x) \leqslant \max_{x\in\partial\Omega} w^+(x).
\end{align*}
于是，我们得到
\begin{align*}
\max_{x\in\overline{\Omega}} u(x)\leqslant \max_{x\in\partial\Omega} w(x) + \varepsilon d^2 \leqslant \max_{x\in\partial\Omega} w^+(x) + \varepsilon d^2 \leqslant \max_{x\in\partial\Omega} u^+(x) + \varepsilon d^2.
\end{align*}
令 $\varepsilon\to 0$，则得到所要证明的结论。至此我们完成定理的证明.
\end{proof}

\begin{theorem}[Hopf(霍夫)引理]\label{theorem:Hopf引理}
假设 $B_R$ 是 $\mathbb{R}^n$ $(n\geqslant 2)$ 上的一个以 $R$ 为半径的球，在 $B_R$ 上 $c(x)\geqslant 0$ 且有界。如果 $u\in C^2(B_R)\cap C^1(\overline{B}_R)$ 满足
\begin{enumerate}[(1)]
\item $\mathcal{L}u = -\Delta u + c(x)u \leqslant 0,\ x\in B_R$；
\item 存在 $x_0\in\partial B_R$ 使得 $u(x)$ 在 $x_0$ 点达到在 $\overline{B}_R$ 上严格的非负最大值，即 $u(x_0)=\max_{\overline{B}_R}u(x)\geqslant 0$，且当 $x\in B_R$ 时，$u(x)<u(x_0)$，
\end{enumerate}
则
\begin{align}
\left.\frac{\partial u}{\partial \nu}\right|_{x=x_0} > 0, \label{eq:hoph-deriv}
\end{align}
其中 $\nu$ 与 $\partial B_R$ 在 $x_0$ 点的单位外法向量 $\bm{n}$ 的夹角小于 $\frac{\pi}{2}$。
\end{theorem}
\begin{remark}
如果 $c(x)\equiv 0$，则证明中的 $v(x)$ 可取为 Laplace 方程的基本解。另外我们可取 $v(x)=e^{-a|x|^2}-e^{-aR^2}$，其中 $a>0$ 足够大。
\end{remark}
\begin{proof}

根据引理的假设，$\left.\frac{\partial u}{\partial \nu}\right|_{x=x_0}\geqslant 0$ 是显然的，我们需要证明它严格大于 $0$。不妨设 $B_R$ 是以原点为球心，半径为 $R$ 的球。在球壳 $B_R^* = \{x\in B_R \mid \frac{R}{2}<|x|<R\}$ 上考虑辅助函数
\begin{align*}
w(x) = u(x) - u(x_0) + \varepsilon v(x),
\end{align*}
其中 $\varepsilon>0$ 和 $v(x)$ 待定。不妨取 $v(x_0)=0$。此时 $w(x_0)=0$。因而 $w(x)$ 在 $\overline{B}_R^*=\{x\in\mathbb{R}^n \mid \frac{R}{2}\leqslant |x|\leqslant R\}$ 上的非负最大值存在。如果能取到 $\varepsilon>0$ 和 $v(x)$ 使得 $w(x)$ 在 $x_0$ 点达到非负最大值 $w(x_0)$，则
\begin{align*}
0 \leqslant \left.\frac{\partial w}{\partial \nu}\right|_{x=x_0} = \left.\frac{\partial u}{\partial \nu}\right|_{x=x_0} + \varepsilon \left.\frac{\partial v}{\partial \nu}\right|_{x=x_0},
\end{align*}
即
\begin{align*}
\left.\frac{\partial u}{\partial \nu}\right|_{x=x_0} \geqslant -\varepsilon \left.\frac{\partial v}{\partial \nu}\right|_{x=x_0}.
\end{align*}
为了使 $w(x)$ 在球壳 $B_R^*$ 的边界上达到非负最大值，由\refthe{theorem:Lu方程极值原理}我们只需要在球壳 $B_R^*$ 上提条件 $\mathcal{L}w\leqslant 0$。因此，我们只需要构造出函数 $v$，使它满足条件
\begin{align*}
\mathcal{L}v \leqslant 0,\quad x\in B_R^*,
\end{align*}
和
\begin{align*}
\left.\frac{\partial v}{\partial \nu}\right|_{x=x_0} < 0,
\end{align*}
就完成定理的证明。

由于球壳 $B_R^*$ 的对称性我们考虑径向对称函数
\begin{align*}
v(x) = |x|^\alpha - R^\alpha,
\end{align*}
其中 $\alpha<0$ 待定。在球壳 $B_R^*$ 上，计算得到
\begin{align*}
\mathcal{L}v &= -\alpha(\alpha+n-2)|x|^{\alpha-2} + c(x)|x|^\alpha - c(x)R^\alpha \leqslant \left[-\alpha(\alpha+n-2) + CR^2\right]|x|^{\alpha-2},
\end{align*}
其中 $C=\sup_{B_R^*} c(x)$。取 $\alpha<0$ 足够小，我们得到 $\mathcal{L}v\leqslant 0$，从而在球壳 $B_R^*$ 上
\begin{align*}
\mathcal{L}w \leqslant 0.
\end{align*}
所以由\refthe{theorem:Lu方程极值原理}知，$w(x)$ 在球壳 $B_R^*$ 的边界 $\partial B_R^*$ 达到非负最大值。在球壳的内球面 $\partial B_R^*\cap B_R = \partial B_{\frac{R}{2}}$ 上，
\begin{align*}
\max_{|x|=\frac{R}{2}} (u(x)-u(x_0)) \text{ 记为 } \beta < 0.
\end{align*}
取 $\varepsilon>0$ 足够小，使得
\begin{align*}
w|_{\partial B_{\frac{R}{2}}} \leqslant \beta + \varepsilon R^\alpha(2^{-\alpha}-1) < 0,
\end{align*}
而在球壳的外球面 $\partial B_R^*\cap\partial B_R = \partial B_R$ 上，$v(x)=0$，显然 $w(x)\leqslant 0$，$w(x_0)=0$，所以 $w(x)$ 在 $x_0$ 点达到非负最大值。由于在 $\partial B_R$ 上 $v(x)=0$，于是
\begin{align*}
\left.\frac{\partial u}{\partial \nu}\right|_{x=x_0} \geqslant -\varepsilon \frac{\partial v}{\partial \nu}\bigg|_{x=x_0} \geqslant -\varepsilon \frac{\partial v}{\partial n}\bigg|_{x=x_0} \cos(\nu,n) = -\varepsilon \alpha R^{\alpha-1}\cos(\nu,n).
\end{align*}
由于 $\alpha$ 是负数，$\nu$ 和 $n$ 的夹角小于 $\frac{\pi}{2}$，我们得到
\begin{align*}
\left.\frac{\partial u}{\partial \nu}\right|_{x=x_0} > 0.
\end{align*}
至此我们完成引理的证明.
\end{proof}

\begin{theorem}[强极值原理]\label{theorem:PDE-强极值原理1}
设 $\Omega$ 是 $\mathbb{R}^n$ 上的有界连通开集，$c(x)\geqslant 0$ 且有界。如果 $u\in C^2(\Omega)\cap C(\overline{\Omega})$ 在 $\Omega$ 上满足 $\mathcal{L}u\leqslant 0$，且 $u(x)$ 在 $\Omega$ 内达到其在 $\overline{\Omega}$ 上的非负最大值，则 $u$ 在 $\Omega$ 上是常数。
\end{theorem}
\begin{proof}
记
\begin{align*}
M = \max_{\overline{\Omega}} u(x).
\end{align*}
考虑集合 $O=\{x\in\Omega \mid u(x)=M\}$。我们只要证明 $O$ 相对于 $\Omega$ 既开又闭，进而由 $\Omega$ 的连通性知，$O$ 是空集或 $\Omega$。又由定理的假设知，$O$ 非空，从而 $O=\Omega$。这就是要证的结论。

由于函数 $u(x)$ 的连续性，$O$ 相对于 $\Omega$ 显然是闭的。现在我们证明 $O$ 相对于 $\Omega$ 是开的。设 $x_0$ 是 $O$ 上任意一点，则存在球 $B(x_0,2r)\subset\Omega$。如果 $x_0$ 不是 $O$ 的内点，则存在 $\tilde{x}\in(\Omega\setminus O)\cap B(x_0,r)$。记
\begin{align*}
d = \operatorname{dist}\{\tilde{x},\overline{O}\} = \min\{|x-\tilde{x}| : x\in\overline{O}\}.
\end{align*}
显然 $d\leqslant r$，因此 $B(\tilde{x},d)\subset B(x_0,2r)\subset\Omega$，
\begin{align*}
u(x) < M, \quad x\in B(\tilde{x},d).
\end{align*}
记 $y_0\in\partial B(\tilde{x},d)\cap O$，则 $u(y_0)=M$。由于 $y_0$ 是函数 $u(x)$ 的极值点，因此
\begin{align*}
\left.\frac{\partial u}{\partial x_i}\right|_{x=y_0} = 0, \quad i=1,2,\cdots,n.
\end{align*}
而在球 $B(\tilde{x},d)$ 上应用\hyperref[theorem:Hopf引理]{Hopf引理}，至少存在一个方向 $\nu$ 使得
\begin{align*}
\left.\frac{\partial u}{\partial \nu}\right|_{x=y_0} > 0.
\end{align*}
这就导致矛盾。因而我们证明了 $x_0$ 是 $O$ 的内点。从而 $O$ 相对于 $\Omega$ 是开的。定理证毕.
\end{proof}

\subsection{最大模估计}

\begin{theorem}\label{theorem:Dirichlet问题的解的最大模估计}
设$\Omega$ 是 $\mathbb{R}^n$ 上的有界开集,$u\in C^2(\Omega)\cap C(\overline{\Omega})$ 是 Dirichlet 问题(第一边值问题)
\begin{align}
\begin{cases}
-\Delta u = f(x), & x\in\Omega, \\
u|_{\partial\Omega} = g.
\end{cases} \label{eq:dirichlet-max}
\end{align}
的解，则
\begin{align*}
\max_{\overline{\Omega}} |u(x)| \leqslant G + CF, 
\end{align*}
其中 $G=\max_{\partial\Omega}|g(x)|$，$F=\sup_{\Omega}|f(x)|$，$C$ 是一个仅依赖于维数 $n$ 以及 $\Omega$ 的直径 $d=\sup_{x,y\in\Omega}|x-y|$ 的常数。
\end{theorem}

\begin{proof}
不妨假设 $\Omega$ 包含原点 $x=0$。令 $w(x)=u(x)-z(x)$，其中
\begin{align*}
z(x) = G + \frac{F}{2n}(d^2 - |x|^2).
\end{align*}
容易验证
\begin{gather*}
-\Delta w = f(x) - F \leqslant 0, \quad x\in\Omega, \\
w|_{\partial\Omega} \leqslant g - G \leqslant 0.
\end{gather*}
由\refthe{theorem:Lu方程极值原理}，在 $\Omega$ 上 $w(x)\leqslant 0$。从而
\begin{align*}
u(x) \leqslant z(x) \leqslant G + \frac{d^2}{2n}F, \quad x\in\overline{\Omega}.
\end{align*}
同理考虑 $-u(x)$ 满足的问题得到
\begin{align*}
u(x) \geqslant -z(x) \geqslant -G - \frac{d^2}{2n}F, \quad x\in\overline{\Omega}.
\end{align*}
因此
\begin{align*}
|u(x)| \leqslant z(x) \leqslant G + \frac{d^2}{2n}F, \quad x\in\overline{\Omega}.
\end{align*}
两边取上确界，定理即得证.
\end{proof}

\begin{definition}
设$\Omega\subset \mathbb{R}^n$,在 $\Omega$ 上考虑边值问题
\begin{align}
\begin{cases}
-\Delta u + c(x)u = f(x), & x\in\Omega, \\
\left[\frac{\partial u}{\partial n} + \alpha(x)u\right]\bigg|_{\partial\Omega} = g(x),
\end{cases} \label{eq:third-boundary}
\end{align}
其中 $n$ 是 $\partial\Omega$ 的单位外法向量。如果 $\alpha(x)\equiv 0$，则问题 \eqref{eq:third-boundary} 称为 \textbf{Neumann 问题}或\textbf{第二边值问题}；如果 $\alpha(x)>0$，则问题 \eqref{eq:third-boundary} 称为\textbf{第三边值问题}。
\end{definition}
\begin{remark}
注意到当 $\alpha(x)\equiv 0$，$c(x)\equiv 0$ 时，齐次 Neumann 问题 $(f(x)\equiv 0, g(x)\equiv 0)$ 有非零解 $u(x)\equiv 1$，因此最大模估计在此情形不成立。但对于第三边值问题，利用极值原理我们可以得到下面的最大模估计。
\end{remark}

\begin{theorem}\label{theorem:边值问题最大模估计11}
设$\Omega$ 是 $\mathbb{R}^n$ 上的有界开集,$c(x)\geqslant 0$，$\alpha(x)\geqslant \alpha_0>0$。如果 $u\in C^2(\Omega)\cap C^1(\overline{\Omega})$ 是问题 \eqref{eq:third-boundary} 的解,则
\begin{align*}
\max_{\overline{\Omega}} |u(x)| \leqslant C(G+F), 
\end{align*}
其中 $G=\max_{\partial\Omega}|g(x)|$，$F=\sup_{\Omega}|f(x)|$，$C$ 是仅依赖于维数 $n$，$\alpha_0$ 和 $\Omega$ 的直径 $d$ 的常数。
\end{theorem}
\begin{remark}
实际上，上述最大模估计蕴涵着第一边值问题 \eqref{eq:dirichlet-max} 和第三边值问题 \eqref{eq:third-boundary} 的解的惟一性和稳定性.
\end{remark}
\begin{proof}
不妨假设 $\Omega$ 包含原点 $x=0$。令 $w(x)=u(x)-z(x)$，其中
\begin{align*}
z(x) = \frac{G}{\alpha_0} + \frac{F}{2n}\left( \frac{1+d^2}{\alpha_0} + d^2 - |x|^2 \right).
\end{align*}
容易验证，在 $\Omega$ 上，
\begin{align*}
-\Delta z + c(x)z \geqslant F;
\end{align*}
在 $\partial\Omega$ 上，
\begin{align*}
\frac{\partial z}{\partial n} + \alpha(x)z
&= \alpha(x)\frac{G}{\alpha_0} + \frac{F}{2n}\left[ -2x\cdot n + \alpha(x)\left( \frac{1+d^2}{\alpha_0} + d^2 - |x|^2 \right) \right] \\
&\geqslant G + \frac{F}{2n}(-|x|^2 - 1 + 1 + d^2) \geqslant G.
\end{align*}
因此，当 $x\in\Omega$ 时，
\begin{align*}
-\Delta w(x) + c(x)w(x) \leqslant f(x) - F \leqslant 0;
\end{align*}
当 $x\in\partial\Omega$ 时，
\begin{align*}
\frac{\partial w}{\partial n} + \alpha(x)w(x) \leqslant g(x) - G \leqslant 0.
\end{align*}
由\refthe{theorem:Lu方程极值原理}，我们知道 $w(x)$ 的正的最大值一定在边界 $\partial\Omega$ 上达到。设在点 $x_0\in\partial\Omega$ 处达到最大值。由于 $n$ 是 $\partial\Omega$ 的单位外法向量，于是 $\left.\frac{\partial w}{\partial n}\right|_{x=x_0}\geqslant 0$，从而
\begin{align*}
\left.\frac{\partial w}{\partial n}\right|_{x=x_0} + \alpha(x_0)w(x_0) \geqslant \alpha(x_0)w(x_0) > 0.
\end{align*}
这与上式矛盾。这说明当 $x\in\overline{\Omega}$ 时，$w(x)\leqslant 0$。所以
\begin{align*}
u(x) \leqslant z(x) \leqslant C(G+F), \quad x\in\overline{\Omega},
\end{align*}
其中
\begin{align*}
C = \max\left\{ \frac{1}{\alpha_0},\ \frac{1}{2n}\left( \frac{1+d^2}{\alpha_0} + d^2 \right) \right\}.
\end{align*}
考虑 $w(x)=u(x)+z(x)$ 我们得到另一方向的不等式，因此
\begin{align*}
|u(x)| \leqslant z(x) \leqslant C(G+F), \quad x\in\overline{\Omega}.
\end{align*}
两边取上确界，定理即证.
\end{proof}

\begin{theorem}
设$\Omega$ 是 $\mathbb{R}^n$ 上的有界开集,$u_i\in C^2(\Omega)\cap C^1(\overline{\Omega})$ $(i=1,2)$ 满足第三边值问题
\begin{align*}
\begin{cases}
-\Delta u_i + c_i(x)u_i = f_i(x), & x\in\Omega, \\
\left[\frac{\partial u_i}{\partial n} + \alpha_i(x)u_i\right]_{\partial\Omega} = g_i(x),
\end{cases}
\end{align*}
其中 $n$ 是 $\partial\Omega$ 的单位外法向量。如果 $c_i(x)\geqslant 0$ 且有界，$\alpha_i(x)\geqslant \alpha_0>0$，则估计
\begin{align}
\max_{\overline{\Omega}} |u_1-u_2| &\leqslant C\big(\max_{\partial\Omega}|g_1-g_2| + \sup_{\Omega}|f_1-f_2|  + \max_{\partial\Omega}|\alpha_1-\alpha_2| + \sup_{\Omega}|c_1-c_2|\big) \label{eq:stability-third}
\end{align}
成立，其中 $C$ 是仅依赖于维数 $n$，$\alpha_0$，$\Omega$ 的直径 $d$ 和 $G_1,G_2,F_1,F_2$ 的常数，这里 $$G_i=\sup_{x\in\partial\Omega}|g_i(x)|,\quad F_i=\sup_{\Omega}|f_i(x)|，\quad i=1,2.$$
\end{theorem}
\begin{proof}
由\refthe{theorem:边值问题最大模估计11} 我们有
\begin{align}
\max_{\overline{\Omega}} |u_i| \leqslant C_1(G_i+F_i), \quad i=1,2. \label{eq:ui-bound}
\end{align}
设 $w=u_1-u_2$，则 $w$ 满足边值问题
\begin{align*}
\begin{cases}
-\Delta w + c_1(x)w = f_1-f_2 + (c_2-c_1)u_2, & x\in\Omega, \\
\left[\frac{\partial w}{\partial n} + \alpha_1(x)w\right]_{\partial\Omega} = g_1-g_2 + (\alpha_2-\alpha_1)u_2.
\end{cases}
\end{align*}
于是，由\refthe{theorem:边值问题最大模估计11} 有
\begin{align*}
\max_{\overline{\Omega}} |w| &\leqslant C_1\big(\max_{\partial\Omega}|g_1-g_2| + \max_{\partial\Omega}|(\alpha_2-\alpha_1)u_2| + \sup_{\Omega}|f_1-f_2| + \sup_{\Omega}|(c_1-c_2)u_2|\big) \\
&\leqslant C_1\big(\max_{\partial\Omega}|g_1-g_2| + \max_{\partial\Omega}|\alpha_2-\alpha_1|\max_{\overline{\Omega}}|u_2|
+ \sup_{\Omega}|f_1-f_2| + \sup_{\Omega}|c_1-c_2|\max_{\overline{\Omega}}|u_2|\big).
\end{align*}
由不等式 \eqref{eq:ui-bound} 我们得到估计 \eqref{eq:stability-third}.
\end{proof}

















\end{document}